\chapter{Arranque del rotor en autorrotación y métodos para inducirlo}\label{cap:arranque}
\usepgfplotslibrary{groupplots}

Con el paso de diseño $\theta=\ang{-3}$ el rotor arranca solo, sin trampas, con un margen de
unos \ang{3} de paso respecto del límite de arranque. Con la apertura de las palas incluida,
llega al \SI{90}{\percent} de sus rpm finales en \SI{1.6}{s}, tras caer \SI{22}{m}
(configuración v2: rotor de $R=\SI{0.25}{m}$, sin drogue, $m_2=\SI{0.49}{kg}$). El
arranque no es un riesgo de vuelo con este paso, pero sí consume altura. Una falla de
arranque deja al CanSat cayendo a unos \SI{14}{m/s} (rotor parado, que frena como un
disco), por encima del límite de C4. Este capítulo
explica la física del arranque desde $\Omega=0$, delimita la región de parámetros en la que el
rotor arranca solo, simula el transitorio con la apertura de las palas y cuantifica once
métodos para acelerar o asegurar el arranque. Todos los números salen del modelo BEMT del
capítulo~\ref{cap:bemt} (\texttt{analysis/bemt\_modelo.py}, importado sin cambios) mediante
el script \texttt{analysis/arranque\_estudios.py}. Son predicciones de modelo; ningún ensayo
de arranque se hizo todavía.

\begin{resultado}[Resultados principales del capítulo]
\begin{itemize}
  \item Con $\Omega=0$ la pala ve el flujo a $\alpha=\ang{90}+\theta$, en pérdida profunda. El
        par de arranque nace de inclinar la fuerza normal ($C_D\approx1{,}17$) hacia adelante:
        $Q_0=\SI{30}{mN.m}$ a \SI{13.4}{m/s} con \ang{-3}, unas 40 veces el rozamiento de
        los rodamientos (\SI{0.73}{mN.m}).
  \item Límite de arranque: con rozamiento nominal el rotor arranca solo si
        $\theta_{0,75}\lesssim\ang{-0.1}$ a \SI{13.4}{m/s} y $\lesssim\ang{-0.3}$ a
        \SI{6.36}{m/s}. Con $\theta\ge\ang{0}$ queda trabado: con \ang{+1} hay un equilibrio
        estable lento (girando hacia atrás a $\approx\SI{3}{rad/s}$) y la caja cae a
        \SI{14.1}{m/s}.
  \item Transitorio base ($\theta=\ang{-3}$, apertura desde \ang{-89}, rozamiento nominal):
        apertura en \SI{0.18}{s}, $t_{90}=\SI{1.55}{s}$, $z_{90}=\SI{21.7}{m}$. La apertura
        cuesta solo \SI{0.06}{s}. El sobrepaso de $\Omega$ llega a \SI{162.6}{rad/s}
        (\num{1553}~rpm), un \SI{26}{\percent} sobre el régimen.
  \item Los métodos más rentables: acoplamiento pasivo paso--batimiento $\delta_3=\ang{30}$
        ($t_{90}=\SI{0.76}{s}$, $z_{90}=\SI{10.8}{m}$, sin masa añadida), paso colectivo variable
        pasivo de \ang{-8} a \ang{-3} ($\SI{0.95}{s}$, \SI{13.4}{m}) y paso fijo \ang{-4}
        ($\SI{1.29}{s}$, \SI{18.0}{m}, costo cero). El resorte de pre-rotación funciona
        ($\Omega_0=\SI{40}{rad/s}$: \SI{0.61}{s}, \SI{0.7}{J}) pero hace girar el cuerpo
        en sentido contrario a $\approx\SI{76}{rad/s}$. Motor, Savonius, liberación más
        temprana y menos rozamiento aportan menos de \SI{0.2}{s}.
  \item \textbf{Recomendación}: mantener paso fijo entre \ang{-3} y \ang{-4} con rodamientos
        limpios (sin cambiar a sellados de contacto); no añadir mecanismos de arranque para el
        primer vuelo. Si se necesita recuperar altura, la opción con mejor relación
        ganancia/riesgo es la bisagra con $\delta_3\approx\ang{30}$, sujeta a la validación
        en banco de la sección~\ref{sec:arranque-ensayo}.
\end{itemize}
\end{resultado}

%=======================================================================================
\section{Física del arranque}\label{sec:arranque-fisica}

El arranque es un paso por la pérdida profunda: la pala parte a $\alpha\approx\ang{87}$ y
solo recupera flujo adherido cerca de las rpm finales. Mientras tanto el par es positivo
porque la fuerza normal de placa, inclinada por el paso negativo, tiene componente hacia
adelante.

\subsection{Triángulo de velocidades con el rotor parado}

En el anillo de radio $r$ la velocidad axial a través del disco es $u_p=V(1-a)$ y la
tangencial $u_t=\Omega r$. El ángulo de entrada y el de ataque son
\begin{equation}
  \phi=\arctan\frac{V(1-a)}{\Omega r},\qquad \alpha=\phi+\theta(r),
  \label{eq:arranque-alfa}
\end{equation}
y la componente de la fuerza de la sección en el plano del disco, positiva hacia adelante
(par motor), es
\begin{equation}
  dQ=\tfrac12\rho U^2 c\,B\,r\,\bigl(C_l\sin\phi-C_d\cos\phi\bigr)\,dr .
  \label{eq:arranque-dq}
\end{equation}
Con $\Omega=0$ resulta $\phi=\ang{90}$, $\alpha=\ang{90}+\theta$ y
$C_l\sin\phi-C_d\cos\phi=C_l(\alpha)$. En pérdida profunda la fuerza es casi normal a la
placa, $C_l\approx C_N\cos\alpha$, de modo que
\begin{equation}
  Q_0\approx\tfrac12\rho V_\text{ef}^2\,c\,B\,C_N\sin(-\theta)\,\frac{R^2-e^2}{2},
  \label{eq:arranque-q0}
\end{equation}
proporcional a $-\theta$ y a $V^2$. El signo de $\theta$ decide el sentido: con
$\theta<0$ el borde de ataque mira hacia arriba-adelante y la fuerza normal empuja la pala
hacia adelante (figura~\ref{fig:arranque-triangulo}); con $\theta>0$ la empuja hacia atrás.
La placa curvada no es simétrica: a \ang{90} su $C_l$ residual es
$0{,}06$ en la polar extendida, que es el término que vuelve $Q_0$ ligeramente positivo a
$\theta=0$ pero menor que el rozamiento.

\begin{figure}[H]
\centering
\begin{tikzpicture}[scale=1.0]
  % pala (placa curvada) con paso negativo
  \begin{scope}[rotate=-8]
    \draw[rot, line width=1.6pt] (-1.8,0) .. controls (-0.6,0.28) and (0.6,0.28) .. (1.8,0);
  \end{scope}
  \node[crot, font=\small] at (2.5,-0.55) {pala, $\theta=\ang{-3}$ (exagerado)};
  \draw[eje] (-3.2,0) -- (3.4,0) node[right, font=\small] {plano del disco};
  % flujo vertical
  \foreach \x in {-1.2,0,1.2}
    \draw[flecha, cfijo, thick] (\x,-2.4) -- (\x,-0.7);
  \node[cfijo, font=\small, anchor=west] at (1.35,-1.9) {$V(1-a)$, $\phi=\ang{90}$};
  % fuerza normal
  \draw[-{Stealth[length=3mm]}, cfreno, very thick] (0,0.15) -- ($(0,0.15)+(98:2.2)$)
     node[above, font=\small] {$dN$ (normal a la placa)};
  \draw[-{Stealth[length=3mm]}, cfreno, thick, dashed] (0,0.15) -- (-0.31,0.15)
     node[pos=1, above left=6pt, font=\small] {$dN\sin(-\theta)$: par};
  \draw[ref] (0,0.15) -- (0,2.4);
  \draw[cgris] (0,1.6) arc (90:98:1.45);
  \node[font=\scriptsize] at (0.35,1.95) {$-\theta$};
  \draw[flecha, crot, thick] (-2.9,1.4) -- (-2.0,1.4) node[midway, above, font=\small] {giro};
  \node[font=\small, align=left, anchor=west] at (-3.3,-2.6)
     {$\alpha=\phi+\theta\approx\ang{87}$\\ $C_N\approx1{,}17$ (alargamiento finito)};
\end{tikzpicture}
\caption{Origen del par de arranque con el rotor parado. El flujo llega normal al disco; la
fuerza normal de la placa, inclinada $-\theta$ por el paso negativo, tiene una componente
hacia adelante. El giro se dibuja hacia la izquierda: el borde de ataque está a la izquierda.}
\label{fig:arranque-triangulo}
\end{figure}

\subsection{Recorrido por la pérdida profunda}

Al crecer $\Omega$, $\phi$ baja y $\alpha$ recorre la pérdida profunda desde \ang{87}
hasta el flujo adherido (figura~\ref{fig:arranque-alfa}). A \SI{13.4}{m/s} la sección de
$0{,}75R$ pasa por \ang{45} a \SI{60}{rad/s} y por \ang{17} a \SI{150}{rad/s}; solo por
encima de \SI{170}{rad/s} se adhiere ($\alpha\approx\ang{5}$). En toda la travesía $C_l$
sube de 0,06 a 0,9 mientras $C_d$ baja de 1,17 a 0,17, y el término
$C_l\sin\phi-C_d\cos\phi$ se mantiene positivo. Las secciones exteriores (mayor $\Omega r$)
se adhieren antes; la raíz sigue en pérdida a rpm de régimen. Por eso el par crece con
$\Omega$ en vez de caer, y la curva $Q(\Omega)$ con $\theta=\ang{-3}$ es monótona hasta
cerca del equilibrio.

\begin{figure}[H]
\centering
\begin{tikzpicture}
\begin{groupplot}[group style={group size=2 by 1, horizontal sep=1.6cm},
  width=0.48\textwidth, height=5.6cm, grid=major, xlabel={$\Omega$ [rad/s]},
  legend style={font=\scriptsize}, xmin=0, xmax=240]
\nextgroupplot[ylabel={$\alpha$ [\si{\degree}]}, ymin=0, ymax=95, legend pos=south west,
  title={\footnotesize $\theta=\ang{-3}$, $V=\SI{13.4}{m/s}$}]
\addplot[cfijo!50, thick] table[x=Om, y=a30] {analysis/data/arranque_alfa.dat};
\addlegendentry{$0{,}3R$}
\addplot[cfijo, thick] table[x=Om, y=a50] {analysis/data/arranque_alfa.dat};
\addlegendentry{$0{,}5R$}
\addplot[crot, very thick] table[x=Om, y=a75] {analysis/data/arranque_alfa.dat};
\addlegendentry{$0{,}75R$}
\addplot[cfreno, thick] table[x=Om, y=a95] {analysis/data/arranque_alfa.dat};
\addlegendentry{$0{,}95R$}
\addplot[ccuerda, thick, dashed] table[x=Om, y=phi75] {analysis/data/arranque_alfa.dat};
\addlegendentry{$\phi$ en $0{,}75R$}
\addplot[black!50, dashed, domain=0:240] {12};
\node[font=\scriptsize, anchor=south west] at (axis cs:100,12) {pérdida $\approx\ang{12}$};
\nextgroupplot[ylabel={coeficientes en $0{,}75R$}, ymin=0, ymax=1.4, legend pos=north east,
  title={\footnotesize polar recorrida}]
\addplot[crot, very thick] table[x=Om, y=cl75] {analysis/data/arranque_alfa.dat};
\addlegendentry{$C_l$}
\addplot[cfijo, very thick] table[x=Om, y=cd75] {analysis/data/arranque_alfa.dat};
\addlegendentry{$C_d$}
\end{groupplot}
\end{tikzpicture}
\caption{Ángulo de ataque y de entrada (izquierda) y coeficientes de la sección de
$0{,}75R$ (derecha) durante el arranque a \SI{13.4}{m/s}. La pala atraviesa toda la pérdida
profunda antes de adherirse.}
\label{fig:arranque-alfa}
\end{figure}

\subsection{Curvas de par y trampas}

Una trampa existe cuando $Q(\Omega)-Q_f$ tiene un cero estable a bajas rpm: el rotor se
queda allí aunque exista la rama rápida. La figura~\ref{fig:arranque-qomega} muestra que con
$\theta=\ang{-3}$ no hay trampa a ninguna $V$ entre 4 y \SI{16}{m/s}: $Q$ es positivo desde
$\Omega=0$ y tiene una sola raíz, la de régimen. Con $\theta=\ang{+1}$ a \SI{13.4}{m/s}
aparecen tres raíces: una estable a $\Omega\approx\SI{-3}{rad/s}$ (el rotor gira despacio
hacia atrás), una inestable a \SI{34}{rad/s} y la estable de régimen a \SI{212}{rad/s}. La
raíz inestable es el umbral: por debajo de ella el rotor vuelve a la trampa. Con $\theta=0$
el par a $\Omega$ pequeño es de \num{0.1} a \SI{1}{mN.m}, del orden del rozamiento: el rotor
queda parado o arranca de forma errática.

\begin{figure}[H]
\centering
\begin{tikzpicture}
\begin{groupplot}[group style={group size=2 by 1, horizontal sep=1.6cm},
  width=0.48\textwidth, height=6cm, grid=major, xlabel={$\Omega$ [rad/s]},
  legend style={font=\scriptsize}]
\nextgroupplot[ylabel={$Q$ [mN·m]}, xmin=0, xmax=400, ymin=-40, ymax=240,
  title={\footnotesize $\theta=\ang{-3}$, varias $V$ [m/s]}, legend pos=north west, legend columns=2]
\addplot[cfijo!40, thick] table[x=Om, y=V4] {analysis/data/arranque_qomega.dat};
\addlegendentry{4}
\addplot[cfijo, thick] table[x=Om, y=V694] {analysis/data/arranque_qomega.dat};
\addlegendentry{6,36}
\addplot[ccont, thick] table[x=Om, y=V10] {analysis/data/arranque_qomega.dat};
\addlegendentry{10}
\addplot[crot, very thick] table[x=Om, y=V134] {analysis/data/arranque_qomega.dat};
\addlegendentry{13,4}
\addplot[cfreno, thick] table[x=Om, y=V16] {analysis/data/arranque_qomega.dat};
\addlegendentry{16}
\addplot[black, thin, domain=0:400] {0};
\nextgroupplot[ylabel={$Q$ [mN·m]}, xmin=-40, xmax=120, ymin=-10, ymax=60,
  title={\footnotesize zona de arranque, $V=\SI{13.4}{m/s}$}, legend pos=north west]
\addplot[crot, very thick] table[x=Om, y=V134] {analysis/data/arranque_qomega.dat};
\addlegendentry{$\theta=\ang{-3}$}
\addplot[crot, thick, dashed] table[x=Om, y=tm3_2D] {analysis/data/arranque_qomega.dat};
\addlegendentry{$\theta=\ang{-3}$, polar 2D}
\addplot[cgris, thick] table[x=Om, y=t0] {analysis/data/arranque_qomega.dat};
\addlegendentry{$\theta=\ang{0}$}
\addplot[cfreno, very thick] table[x=Om, y=t1] {analysis/data/arranque_qomega.dat};
\addlegendentry{$\theta=\ang{+1}$}
\addplot[black, thin, domain=-40:120] {0};
\addplot[black!60, dashed, domain=-40:120] {0.73};
\addplot[only marks, mark=*, mark size=2pt, ccuerda] coordinates {(-3,0)};
\addplot[only marks, mark=o, mark size=2.5pt, cfreno, thick] coordinates {(41,0)};
\node[font=\scriptsize, anchor=north] at (axis cs:-5,-1) {trampa};
\node[font=\scriptsize, anchor=north] at (axis cs:48,-1) {umbral};
\end{groupplot}
\end{tikzpicture}
\caption{Curvas $Q(\Omega)$. Izquierda: paso de diseño a varias velocidades. Derecha: zona de
arranque a \SI{13.4}{m/s}; la línea de trazos es $Q_f=\SI{0.73}{mN.m}$. Con \ang{+1} el
cero estable a $\Omega<0$ es la trampa y el cero inestable a \SI{34}{rad/s} el umbral.}
\label{fig:arranque-qomega}
\end{figure}

\begin{nota}[Validez a $\Omega\le0$]
La rama $\Omega<0$ (pala girando hacia atrás, flujo entrando por el borde de fuga) se evaluó
con la misma polar extendida $\pm\ang{180}$ y el mismo cierre de cantidad de movimiento,
que con el rotor casi parado opera en la rama extrapolada $a>1$
(sec.~\ref{sec:bemt-cierre}). La existencia de la trampa con $\theta\ge\ang{+1}$ es robusta,
porque depende solo del signo de $Q_0$; su posición exacta no lo es.
\end{nota}

\subsection{Mapa de arranque}

El rotor arranca solo si $Q(0)>Q_f$ y $Q(\Omega)-Q_f$ cambia de signo una sola vez entre
$\Omega=0$ y la rama de régimen. Ese criterio se evaluó en una malla de 20 valores de
$\Omega$ entre 0 y \SI{420}{rad/s}, y se buscó por bisección (resolución \ang{0.1}) el paso
crítico $\theta_c(V)$: el rotor arranca para $\theta_{0,75}<\theta_c$. La
figura~\ref{fig:arranque-mapa} y la tabla~\ref{tab:arranque-mapa} dan el resultado.

\begin{figure}[H]
\centering
\begin{tikzpicture}
\begin{axis}[width=0.8\textwidth, height=6.6cm, grid=major, xlabel={$V$ al soltar [m/s]},
  ylabel={$\theta_c$ (paso en $0{,}75R$) [\si{\degree}]}, xmin=4, xmax=16, ymin=-10, ymax=2,
  legend style={font=\scriptsize, at={(0.5,-0.22)}, anchor=north}, legend columns=4]
\addplot[name path=nom, crot, very thick, mark=*, mark size=1.3pt] table[x=V, y=f073] {analysis/data/arranque_mapa.dat};
\addlegendentry{$Q_f=0{,}73$ (nominal)}
\addplot[crot!50, thick, dashed] table[x=V, y=f034] {analysis/data/arranque_mapa.dat};
\addlegendentry{$Q_f=0{,}34$}
\addplot[cfreno, thick, dashed] table[x=V, y=f17] {analysis/data/arranque_mapa.dat};
\addlegendentry{$Q_f=1{,}7$}
\addplot[cfreno, thick, dotted] table[x=V, y=f5] {analysis/data/arranque_mapa.dat};
\addlegendentry{$Q_f=5$ (sucio/sellos)}
\addplot[ccuerda, thick, mark=triangle*, mark size=1.5pt] table[x=V, y=tw10] {analysis/data/arranque_mapa.dat};
\addlegendentry{torsión $+\ang{10}$}
\addplot[ccuerda!60!black, thick, mark=square*, mark size=1.2pt] table[x=V, y=tw20] {analysis/data/arranque_mapa.dat};
\addlegendentry{torsión $+\ang{20}$}
\addplot[cfijo, thick, densely dashdotted] table[x=V, y=p2D] {analysis/data/arranque_mapa.dat};
\addlegendentry{polar 2D}
\addplot[black, thick, domain=4:16] {-3};
\node[font=\scriptsize, anchor=south west] at (axis cs:4.2,-3) {diseño $\theta=\ang{-3}$};
\node[font=\scriptsize, fill=white, inner sep=1pt] at (axis cs:12,-5) {ARRANCA SOLO};
\node[font=\scriptsize, fill=white, inner sep=1pt] at (axis cs:12,1.5) {NO ARRANCA};
\end{axis}
\end{tikzpicture}
\caption{Mapa de arranque: paso crítico por encima del cual el rotor no arranca solo, en
función de la velocidad de liberación, del rozamiento $Q_f$ [mN·m], de la torsión lineal y
de la polar. Debajo de cada curva el rotor arranca sin ayuda.}
\label{fig:arranque-mapa}
\end{figure}

\begin{table}[H]
\centering
\caption{Paso crítico de arranque $\theta_c$ [\si{\degree}] (en $0{,}75R$) y margen del
diseño ($\theta=\ang{-3}$). Resolución de la bisección: \ang{0.1}.}
\label{tab:arranque-mapa}
\begin{tabular}{lrrrrrrr}
\toprule
$V$ [m/s] & 4 & 5 & 6,36 & 9 & 11 & 13,4 & 16 \\
\midrule
$Q_f=0{,}73$ mN·m (nominal) & $-0{,}9$ & $-0{,}5$ & $-0{,}3$ & $-0{,}2$ & $-0{,}1$ & $-0{,}1$ & $-0{,}1$ \\
$Q_f=0{,}34$ mN·m (aceite) & $-0{,}4$ & $-0{,}2$ & $-0{,}1$ & $-0{,}1$ & $-0{,}1$ & $0{,}0$ & $0{,}0$ \\
$Q_f=1{,}7$ mN·m (grasa fría) & $-2{,}0$ & $-1{,}2$ & $-0{,}8$ & $-0{,}4$ & $-0{,}2$ & $-0{,}2$ & $-0{,}1$ \\
$Q_f=5$ mN·m (sellos, suciedad) & $-5{,}6$ & $-3{,}6$ & $-2{,}2$ & $-1{,}1$ & $-0{,}8$ & $-0{,}5$ & $-0{,}3$ \\
torsión $+\ang{10}$, $Q_f$ nominal & $-0{,}1$ & $+0{,}2$ & $+0{,}3$ & $+0{,}4$ & $+0{,}6$ & $+0{,}6$ & $+0{,}6$ \\
torsión $+\ang{20}$, $Q_f$ nominal & $+0{,}6$ & $+0{,}8$ & $+1{,}0$ & $+1{,}1$ & $+1{,}2$ & $+1{,}2$ & $+1{,}2$ \\
polar 2D ($C_{D,\max}=1{,}98$) & $-0{,}5$ & $-0{,}3$ & $-0{,}2$ & $-0{,}1$ & $-0{,}1$ & $-0{,}1$ & $0{,}0$ \\
\midrule
margen del diseño, nominal & 2,1 & 2,5 & 2,7 & 2,8 & 2,9 & 2,9 & 2,9 \\
\bottomrule
\end{tabular}
\end{table}

Tres conclusiones. Primera, a la velocidad de liberación (\SI{13.4}{m/s}) el límite queda
prácticamente en \ang{0}, lo que confirma la regla del capítulo~\ref{cap:dim}; el diseño
tiene \ang{2.9} de margen, más que la tolerancia de fabricación de $\pm\ang{1}$. Segunda, el
rozamiento solo importa a baja velocidad o si se multiplica: con \SI{5}{mN.m} (sellos de
contacto, suciedad, grasa fría) y una liberación lenta el paso de diseño quedaría justo en el
límite (a \SI{4}{m/s} ya no arrancaría). Tercera, la torsión con raíz más negativa
($\theta_{tw}>0$ en la convención
$\theta(r)=\theta_{0,75}+\theta_{tw}(r/R-0{,}75)$) mueve el límite hacia pasos positivos:
\ang{0.6} con $\theta_{tw}=\ang{+10}$ y \ang{1.2} con \ang{+20}, porque la raíz, que aporta
el par a $\Omega$ bajo, trabaja con más paso negativo. Es el mismo sentido que en el
capítulo~\ref{cap:bemt}.

\paragraph{Sensibilidad a la polar en pérdida.} Con la polar 2D ($C_{D,\max}=1{,}98$ en vez
de 1,17), $Q_0$ sube de \num{30} a \SI{44}{mN.m} a \SI{13.4}{m/s} (\SI{+47}{\percent}),
pero el paso crítico casi no cambia: el límite lo fija el signo de la fuerza en el plano, no
su módulo. La polar en pérdida afecta el tiempo de arranque (\SI{-0.18}{s},
sec.~\ref{sec:arranque-dinamica}), no la existencia del arranque.

%=======================================================================================
\section{Dinámica temporal del arranque}\label{sec:arranque-dinamica}

Con la apertura de las palas incluida, el rotor llega al \SI{90}{\percent} de sus rpm en
\SI{1.55}{s} y \SI{21.7}{m}, solo \SI{0.06}{s} más que el cálculo cuasiestacionario del
capítulo~\ref{cap:bemt}.

\subsection{Modelo}

El estado es $(\Omega, V, \beta, \dot\beta, z)$. El rotor conserva su momento cinético en
el eje:
\begin{equation}
  \frac{d}{dt}\bigl[I_z(\beta)\,\Omega\bigr]=f(\beta)\,Q(\theta,\Omega,V)-Q_f\operatorname{sgn}\Omega+Q_\text{ext}(t),
  \qquad
  m\dot V=W-f(\beta)\,T(\theta,\Omega,V)-\tfrac12\rho V^2C_{D,c}A_c ,
  \label{eq:arranque-din}
\end{equation}
con $T$ y $Q$ interpolados en tablas BEMT $(\theta,\Omega,V)$ de
$10\times31\times13$ puntos ($\theta=\ang{-10}$ a \ang{+1}, $\Omega=-40$ a \SI{350}{rad/s},
$V=2$ a \SI{17}{m/s}). El factor $f=\cos^2\beta$ para $\beta<0$ (pala colgando) reduce
fuerza y brazo; $f=1$ con la pala abierta. La inercia polar depende de la apertura:
\begin{equation}
  I_z(\beta)=I_\text{cubo}+B\bigl(I_b\cos^2\beta+2S_be\cos\beta+m_be^2\bigr),
  \label{eq:arranque-iz}
\end{equation}
con $I_b=\SI{9.74e-5}{kg.m^2}$, $m_b=\SI{9.5}{g}$, $S_b=m_bx_g$ ($x_g\approx\SI{92}{mm}$,
estimado por escala de la pala de $R=\SI{0.20}{m}$), $e=\SI{44}{mm}$ e $I_\text{cubo}$ del
rotor anterior ($\SI{8.6e-5}{kg.m^2}$, cubo y herrajes sin cambio). Resulta
$I_z(0)=\SI{8.6e-4}{kg.m^2}$ e $I_z(\ang{-89})=\SI{1.65e-4}{kg.m^2}$: con las palas colgando
el rotor tiene el \SI{19}{\percent} de su inercia final.

La apertura sigue el batimiento simplificado de la ecuación~\eqref{eq:estructura-apertura}
del capítulo de estructura, con fuerza de flujo cruzado ($C_{dc}=1{,}2$), resorte de
\SI{0.005}{N.m}, peso y momento centrífugo $-\Omega^2\sin\beta\,(I_b\cos\beta+S_be)$, desde
$\beta=\ang{-89}$. Cuando $\beta$ alcanza la conicidad cuasiestática
\begin{equation}
  \beta_c=\min\!\left[\ang{15},\;k\,\frac{(T/B)\,0{,}7L_b-S_bg}{I_b\Omega^2}\right],
  \label{eq:arranque-cono}
\end{equation}
se da la pala por abierta y $\beta=\beta_c$ desde entonces; el tope elástico de \ang{+15}
limita la conicidad a bajas rpm. El factor $k=0{,}8$ calibra $\beta_c\approx\ang{4.5}$ en
régimen. El rebote en el tope, el batimiento cíclico y la inercia del flujo inducido se
ignoran. El paso de integración es \SI{2}{ms} (Euler semi-implícito), con $t_{90}$ y
$z_{90}$ medidos sobre el valor de $\Omega$ a \SI{8}{s}.

\subsection{Transitorio de referencia}

La figura~\ref{fig:arranque-trans} muestra el caso base. Las palas se abren en
\SI{0.18}{s} (con el rotor más largo la apertura es algo más lenta que los
\SIrange{127}{148}{ms} del rotor de \SI{0.20}{m}) y quedan contra el tope de \ang{15} hasta
$\approx\SI{1.4}{s}$, porque la fuerza centrífuga es baja. El primer segundo es lento: el
par es de 36 a \SI{61}{mN.m} y la aceleración media de unos \SI{45}{rad/s^2}. Entre 1,2 y
\SI{2}{s}, ya con flujo casi adherido en la punta, el par llega a \SI{212}{mN.m} (a
\SI{1.6}{s}) y $\Omega$ sube hasta un sobrepaso de \SI{162.6}{rad/s} a \SI{1.9}{s}, un
\SI{26}{\percent} sobre el régimen. El régimen es $\Omega=\SI{128.7}{rad/s}$,
$V=\SI{6.36}{m/s}$, igual al del capítulo~\ref{cap:bemt}. Ese sobrepaso (\num{1553}~rpm)
es la velocidad de diseño de las cargas centrífugas (capítulo~\ref{cap:estructura}).

\begin{figure}[H]
\centering
\begin{tikzpicture}
\begin{groupplot}[group style={group size=2 by 1, horizontal sep=1.6cm},
  width=0.48\textwidth, height=5.8cm, grid=major, xlabel={$t$ [s]}, xmin=0, xmax=6,
  legend style={font=\scriptsize}]
\nextgroupplot[ylabel={$\Omega$ [rad/s]}, ymin=-10, ymax=200, legend pos=south east]
\addplot[crot, very thick] table[x=t, y=Om] {analysis/data/arranque_hist_base.dat};
\addlegendentry{base \ang{-3}}
\addplot[crot!50, thick, dashed] table[x=t, y=Om] {analysis/data/arranque_hist_qs.dat};
\addlegendentry{sin apertura}
\addplot[cfijo, thick] table[x=t, y=Om] {analysis/data/arranque_hist_tm5.dat};
\addlegendentry{\ang{-5} fijo}
\addplot[ccont, thick] table[x=t, y=Om] {analysis/data/arranque_hist_var.dat};
\addlegendentry{paso variable}
\addplot[ccuerda, thick] table[x=t, y=Om] {analysis/data/arranque_hist_d3_30.dat};
\addlegendentry{$\delta_3=\ang{30}$}
\addplot[cfreno, thick] table[x=t, y=Om] {analysis/data/arranque_hist_t1.dat};
\addlegendentry{\ang{+1} (trabado)}
\nextgroupplot[ylabel={$V$ [m/s] \;/\; $\beta$ [\si{\degree}]}, ymin=0, ymax=16, legend pos=north east]
\addplot[cfijo, very thick] table[x=t, y=V] {analysis/data/arranque_hist_base.dat};
\addlegendentry{$V$ base}
\addplot[cfreno, thick, dashed] table[x=t, y=V] {analysis/data/arranque_hist_t1.dat};
\addlegendentry{$V$ trabado}
\addplot[crot, thick] table[x=t, y=beta] {analysis/data/arranque_hist_base.dat};
\addlegendentry{$\beta$ base ($\ge0$)}
\end{groupplot}
\end{tikzpicture}
\caption{Arranque desde $\Omega=0$, $V=\SI{13.4}{m/s}$, rotor solo y apertura de palas desde
\ang{-89}. Izquierda: rpm para varios métodos. Derecha: velocidad de caída y conicidad del
caso base (la conicidad se dibuja desde que la pala está abierta).}
\label{fig:arranque-trans}
\end{figure}

\paragraph{Efecto de la apertura.} La apertura cuesta solo \SI{0.06}{s} y \SI{1.6}{m}, por
dos razones: dura poco frente a $t_{90}$, y mientras las palas cuelgan la inercia es la
quinta parte, de modo que el poco par disponible acelera igual. La conservación del momento
cinético tiene otra consecuencia, que es la clave para la pre-rotación: un giro dado
\emph{antes} de abrir las palas se divide por $I_z(0)/I_z(\ang{-89})=5{,}2$ al abrirse.

\paragraph{Soltar el drogue antes de abrir (configuración v2).} Sin el drogue la caja
acelera mientras las palas se abren: en \SI{0.2}{s} gana \SI{1.7}{m/s} y en \SI{0.3}{s},
\SI{2.5}{m/s} (capítulo~\ref{cap:arq}). Con la liberación a \SI{16}{m/s} (casos 13),
$t_{90}$ baja \SI{0.17}{s}: el par crece con $V^2$ y la caída libre previa ayuda al
arranque, aunque el rotor tiene que frenar desde más velocidad (más sobrepaso). Con la
Arquitectura~2 ($m_2=\SI{0.36}{kg}$) el arranque tarda \SI{1.65}{s} y \SI{20.9}{m}
(caso 25), y el régimen baja a \SI{111}{rad/s} y \SI{5.46}{m/s}; la carga más liviana frena
menos, por lo que tarda \SI{0.1}{s} más en llegar al régimen.

\paragraph{Caso trabado.} Con \ang{+1} el rotor gira despacio hacia atrás
($\Omega\to\SI{-2.8}{rad/s}$) y la caja desciende a \SI{14.1}{m/s}, fuera del requisito
C4. Con \ang{0} queda parado ($Q<Q_f$) y el resultado es el mismo.

\subsection{Tiempos y alturas de arranque}

La tabla~\ref{tab:arranque-casos} resume los 26 casos simulados. Todos parten de
$\Omega=0$ y \SI{13.4}{m/s} con apertura de palas, salvo los indicados. La altura $z_{90}$ es
la caída desde la liberación hasta el \SI{90}{\percent} de las rpm finales.

\begin{table}[H]
\centering
\small
\caption{Arranque simulado: tiempo y caída hasta el \SI{90}{\percent} de $\Omega$ final y
régimen a \SI{8}{s}. $Q_f=\SI{0.73}{mN.m}$ salvo indicación.}
\label{tab:arranque-casos}
\begin{tabular}{rlrrrr}
\toprule
\# & Caso & $t_{90}$ [s] & $z_{90}$ [m] & $\Omega_f$ [rad/s] & $V_f$ [m/s] \\
\midrule
1 & Base, $\theta=\ang{-3}$ & 1,55 & 21,7 & 128,7 & 6,36 \\
2 & Base sin apertura (cuasiestacionario, cap.~\ref{cap:bemt}) & 1,49 & 20,0 & 128,7 & 6,36 \\
3 & $\theta=\ang{-2}$ & 2,02 & 28,1 & 122,1 & 6,30 \\
4 & $\theta=\ang{-4}$ & 1,29 & 18,0 & 135,1 & 6,43 \\
5 & $\theta=\ang{-5}$ & 1,12 & 15,7 & 141,2 & 6,51 \\
6 & $\theta=\ang{-6}$ & 1,00 & 14,0 & 147,2 & 6,60 \\
7 & Torsión $+\ang{10}$ ($\theta_{0,75}=\ang{-3}$, raíz \ang{-8.3}) & 1,40 & 19,6 & 128,2 & 6,25 \\
8 & Torsión $+\ang{20}$, $\theta_{0,75}=\ang{-1}$ & 1,99 & 27,6 & 112,7 & 6,19 \\
9 & Pestaña: \ang{-4} equivalentes en la mitad externa (fija) & 1,01 & 14,2 & 147,1 & 6,64 \\
10 & Polar 2D ($C_{D,\max}=1{,}98$) & 1,38 & 17,4 & 129,0 & 6,36 \\
11 & $Q_f=\SI{0.34}{mN.m}$ (aceite) & 1,54 & 21,5 & 128,9 & 6,35 \\
12 & $Q_f=\SI{1.7}{mN.m}$ & 1,58 & 22,0 & 128,3 & 6,36 \\
13 & Liberación a \SI{16}{m/s} & 1,39 & 21,1 & 128,7 & 6,36 \\
14 & Liberación a \SI{10}{m/s} & 1,80 & 22,5 & 128,7 & 6,36 \\
15 & Liberación a \SI{6.36}{m/s} & 2,11 & 23,4 & 128,7 & 6,36 \\
16 & Paso variable pasivo \ang{-8}$\to$\ang{-3} ($\Omega_s=\SI{80}{rad/s}$) & 0,95 & 13,4 & 128,7 & 6,36 \\
17 & $\delta_3=\ang{30}$ & 0,76 & 10,8 & 129,2 & 6,36 \\
18 & $\delta_3=\ang{45}$ & 0,71 & 10,1 & 129,4 & 6,36 \\
19 & Motor \SI{3}{mN.m} durante \SI{1}{s} & 1,48 & 20,3 & 128,7 & 6,36 \\
20 & Resorte, $\Omega_0=\SI{20}{rad/s}$ (palas abiertas) & 0,98 & 13,0 & 128,7 & 6,36 \\
21 & Resorte, $\Omega_0=\SI{40}{rad/s}$ & 0,61 & 7,9 & 128,7 & 6,36 \\
22 & Resorte, $\Omega_0=\SI{80}{rad/s}$ & 0,18 & 2,3 & 128,7 & 6,36 \\
\midrule
23 & $\theta=\ang{+1}$ sin ayuda & --- & --- & $-2{,}8$ & 14,07 \\
24 & $\theta=\ang{0}$ sin ayuda & --- & --- & 0 & 14,08 \\
\midrule
25 & Arq.~2: $m_2=\SI{0.36}{kg}$ & 1,65 & 20,9 & 111,2 & 5,46 \\
26 & Arq.~2: $m_2=\SI{0.36}{kg}$, apertura a \SI{16}{m/s} & 1,48 & 20,3 & 111,2 & 5,46 \\
\bottomrule
\end{tabular}
\end{table}

El valor de referencia del capítulo~\ref{cap:bemt} (\SI{1.5}{s}, \SI{20}{m}) se reproduce
(caso 2). La incertidumbre de la polar en pérdida (casos 1 y 10) vale \SI{0.18}{s} y
\SI{4}{m}; la del rozamiento, \SI{0.04}{s}. Estas cifras dan la escala de lo que un método
de arranque debe superar para que su efecto sea distinguible del error del modelo.

%=======================================================================================
\section{Métodos para inducir el arranque}\label{sec:arranque-metodos}

Los métodos se ordenan en dos familias. Los que aumentan el par aerodinámico a $\Omega$ bajo
(paso, torsión, paso variable, $\delta_3$, pestañas, Savonius, liberación a mayor $V$) y los
que aportan momento cinético o reducen pérdidas (resorte, servo, motor, rozamiento). La
figura~\ref{fig:arranque-tiempos} resume los dos barridos continuos.

\begin{figure}[H]
\centering
\begin{tikzpicture}
\begin{groupplot}[group style={group size=2 by 1, horizontal sep=1.8cm},
  width=0.48\textwidth, height=5.6cm, grid=major, legend style={font=\scriptsize}]
\nextgroupplot[xlabel={paso fijo $\theta$ [\si{\degree}]}, ylabel={$t_{90}$ [s]}, xmin=-10, xmax=0,
  ymin=0, ymax=6, axis y line*=left]
\addplot[crot, very thick, mark=*, mark size=1.4pt] table[x=th, y=t90] {analysis/data/arranque_tpaso.dat};
\addplot[only marks, mark=square*, black, mark size=2pt] coordinates {(-3,2.394)};
\node[font=\scriptsize, anchor=south west] at (axis cs:-3,2.6) {diseño};
\nextgroupplot[xlabel={$\Omega_0$ del resorte [rad/s]}, ylabel={$t_{90}$ [s] \,/\, $z_{90}/10$ [m]},
  xmin=0, xmax=120, ymin=0, ymax=3, legend pos=north east]
\addplot[ccont, very thick, mark=*, mark size=1.4pt] table[x=Om0, y=t90] {analysis/data/arranque_resorte.dat};
\addlegendentry{$t_{90}$}
\addplot[ccuerda, thick, mark=triangle*, mark size=1.5pt] table[x=Om0, y expr=\thisrow{z90}/10] {analysis/data/arranque_resorte.dat};
\addlegendentry{$z_{90}/10$}
\addplot[cfreno, thick, dashed] table[x=Om0, y=E] {analysis/data/arranque_resorte.dat};
\addlegendentry{energía [J]}
\end{groupplot}
\end{tikzpicture}
\caption{Izquierda: tiempo de arranque con paso fijo (con \ang{-0.5} el arranque tarda
\SI{8.9}{s}; a \ang{0} no ocurre). Derecha: efecto de la pre-rotación $\Omega_0$ con palas
abiertas, con $\theta=\ang{-3}$, y energía cinética $\frac12I_z\Omega_0^2$ correspondiente.}
\label{fig:arranque-tiempos}
\end{figure}

\subsection{Paso más negativo}

Cada grado de paso negativo adicional quita unos \SI{0.25}{s} y \SI{3}{m} de arranque, a
costa de rpm: de \ang{-3} a \ang{-4}, $t_{90}$ baja de 1,55 a \SI{1.29}{s} y la caída de 21,7
a \SI{18.0}{m}, mientras $\Omega_f$ sube de 129 a \SI{135}{rad/s} y $V_2$ de 6,36 a
\SI{6.43}{m/s}. A \ang{-6} se ganan \SI{0.55}{s} y \SI{7.6}{m}, con \num{1406}~rpm y
$V_2=\SI{6.60}{m/s}$, todavía dentro de C4 pero con más carga centrífuga y menos margen de
flameo (cap.~\ref{cap:estructura}). El rendimiento decrece: por debajo de \ang{-6} el
arranque solo mejora \SI{0.1}{s} por grado. La curva también muestra lo contrario: hacia
\ang{-1} el arranque se alarga a \SI{3.2}{s} y \SI{44}{m}, y a \ang{-0.5} a \SI{5.0}{s} y
\SI{69}{m}. Ese es el riesgo real de un portapala mal impreso.

\subsection{Torsión y paso variable en la envergadura}

La torsión con raíz más negativa ayuda poco al tiempo de arranque: $\theta_{tw}=\ang{+10}$
con $\theta_{0,75}=\ang{-3}$ (raíz a \ang{-8.3}, punta a \ang{-0.5}) gana \SI{0.15}{s} y
\SI{2.1}{m} sin cambiar el régimen. Su valor está en el margen: desplaza el paso crítico
\ang{0.6} (tabla~\ref{tab:arranque-mapa}). En la placa de carbono de \SI{0.5}{mm}, una
torsión de \ang{10} sobre \SI{206}{mm} exige un molde torsionado; es factible pero añade
una variable de fabricación. Una alternativa equivalente, más sencilla, es un portapala con
\ang{2} a \ang{3} más de paso negativo en la raíz solo en el tramo de unión (paso escalonado),
que aporta en el arranque el mismo efecto del primer orden.

\subsection{Pre-rotación por resorte en el cubo}

Un resorte que dé $\Omega_0$ al soltar acorta el arranque casi linealmente: \SI{-0.025}{s}
por cada rad/s, hasta $\Omega_0\approx\SI{100}{rad/s}$ (figura~\ref{fig:arranque-tiempos}).
Su función depende del paso.

\paragraph{Con paso de diseño.} No hace falta para arrancar, solo para ganar altura:
$\Omega_0=\SI{40}{rad/s}$ da $t_{90}=\SI{0.61}{s}$ y $z_{90}=\SI{7.9}{m}$ (\SI{-13.8}{m}).

\paragraph{Con paso $\ge0$ (trampa).} El resorte debe llevar el rotor por encima del umbral
inestable. Por bisección de la simulación se obtuvo el $\Omega_0$ mínimo
(tabla~\ref{tab:arranque-om0}). Con \ang{+1} se necesitan \num{28} a \SI{37}{rad/s}, algo
menos que el cero de $Q$ (\SI{34}{rad/s} a \SI{13.4}{m/s}) porque la caída se frena mientras
tanto. Con \ang{0} basta $\approx\SI{15}{rad/s}$ (valor acotado por la ventana de simulación
de \SI{4}{s}, por lo lento de la aceleración).

\begin{table}[H]
\centering
\caption{Velocidad inicial mínima $\Omega_0$ [rad/s] (palas abiertas) para salir de la
trampa, y energía del resorte $E=\tfrac12I_z\Omega_0^2$.}
\label{tab:arranque-om0}
\begin{tabular}{lrrrr}
\toprule
$V$ al soltar [m/s] & 6,36 & 10 & 13,4 & 16 \\
\midrule
$\theta=\ang{0}$: $\Omega_0$ & $\approx15$ & $\approx15$ & $\approx15$ & $\approx15$ \\
$\theta=\ang{+1}$: $\Omega_0$ & 27,8 & 31,1 & 34,6 & 37,2 \\
$\theta=\ang{+1}$: $E$ [J] & 0,33 & 0,41 & 0,51 & 0,59 \\
\bottomrule
\end{tabular}
\end{table}

\paragraph{Dimensionamiento.} Se toma $\Omega_0=\SI{40}{rad/s}$ como objetivo (salida de
trampa con \ang{+1} y ganancia de \SI{14}{m} con \ang{-3}): $E=\SI{0.69}{J}$. El resorte debe
actuar \emph{después} de abrir las palas: si actúa antes, el momento cinético se reparte en una
inercia 5,2 veces mayor y la energía necesaria sube a \SI{3.6}{J}. Con un resorte de
torsión de una vuelta útil ($\varphi=2\pi$), $k=2E/\varphi^2=\SI{0.035}{N.m/rad}$ y
$M_\text{max}=k\varphi=\SI{0.22}{N.m}$. Con alambre de acero de resorte y tensión de flexión
admisible \SI{1000}{MPa}~\cite{budynas2015},
\begin{equation}
  d=\left(\frac{32M_\text{max}}{\pi\sigma_\text{adm}}\right)^{1/3}=\SI{1.3}{mm},\qquad
  N=\frac{E_sd^4}{64kD}=\num{11}\text{ espiras}\;(D=\SI{20}{mm}),
  \label{eq:arranque-resorte}
\end{equation}
lo que da unos \SI{0.7}{m} de alambre y \SI{7}{g}, más trinquete y gatillo (\SIrange{3}{5}{g}
estimados).

\begin{nota}[Reacción sobre el cuerpo]
El resorte se apoya en el eje fijo, es decir, en el cuerpo. Sin par externo, el cuerpo recibe
el momento cinético opuesto: con $I_\text{cuerpo}\approx\SI{4.5e-4}{kg.m^2}$ gira a
$-I_z\Omega_0/I_\text{cuerpo}\approx\SI{-76}{rad/s}$ (\num{730}~rpm). Es cien veces el criterio
de video del capítulo~\ref{cap:estabilidad} y el drogue y las aletas tardarían segundos en
amortiguarlo. Esta es la objeción principal a toda pre-rotación con reacción interna (resorte,
servo, motor).
\end{nota}

\subsection{Pre-rotación con el servo o con un motor}

El servo del equipo (SG90, \SI{0.18}{N.m} según su hoja de datos; el MG90S da algo más) pero gira a \SI{60}{\degree} en \SI{0.1}{s}, unos
\SI{10}{rad/s} (valores de catálogo típicos); acoplado directo al cubo con un trinquete no
puede dar más que $\Omega_0\approx\SI{10}{rad/s}$, que vale \SI{0.3}{s} de ganancia. Además
ya está asignado a la traba. Un motor sin escobillas o con núcleo hueco pequeño
(\SIrange{3}{5}{g}) con \SI{3}{mN.m} durante \SI{1}{s} gana solo \SI{0.07}{s} y \SI{1.3}{m}
(caso 19): su par es la sexta parte del aerodinámico. Ambos comparten la reacción sobre el
cuerpo.

\subsection{Aprovechar la apertura de las palas}

La apertura entrega a cada pala un impulso angular alrededor de la bisagra de
$I_b\dot\beta_i\approx\SI{4.7e-3}{N.m.s}$ ($\dot\beta_i\approx\SI{48}{rad/s}$ estimado, la
energía de golpe del cap.~\ref{cap:estructura} escalada al rotor de \SI{0.25}{m}). Si el eje de la bisagra se
gira un ángulo $\delta$ en el plano del disco, la fuerza aerodinámica que abre la pala tiene
una componente tangencial y una fracción $\sin\delta$ de ese impulso pasa al eje del rotor:
\begin{equation}
  \Omega_0\approx\frac{B\,I_b\dot\beta_i\sin\delta}{I_z}\approx\SI{22}{rad/s}\cdot\sin\delta .
  \label{eq:arranque-apert}
\end{equation}
Con $\delta=\ang{20}$, $\Omega_0\approx\SI{7.5}{rad/s}$ y la ganancia es \SI{0.2}{s} (estimación
de orden de magnitud). Lo relevante es que una bisagra girada es justamente un acoplamiento
$\delta_3$: el beneficio grande no viene del impulso de apertura sino del cambio de paso con
la conicidad (sección~\ref{sec:arranque-d3}). A diferencia del resorte, aquí el par sale del
aire y no reacciona sobre el cuerpo.

\subsection{Paso colectivo variable}

Un paso grande al inicio y de régimen después combina el arranque de \ang{-8} con las rpm de
\ang{-3}. Se simuló la ley pasiva
\begin{equation}
  \theta(\Omega)=\theta_r+(\theta_i-\theta_r)\max\!\left[0,\;1-(\Omega/\Omega_s)^2\right],
  \qquad \theta_i=\ang{-8},\;\theta_r=\ang{-3},\;\Omega_s=\SI{80}{rad/s},
  \label{eq:arranque-var}
\end{equation}
que representa un contrapeso centrífugo contra un resorte (fuerza $\propto\Omega^2$). El
arranque baja a \SI{0.95}{s} y \SI{13.4}{m} (\SI{-0.60}{s}, \SI{-8.3}{m}) sin cambiar el
régimen. Su modo de falla es benigno: trabado a \ang{-8} el rotor gira a \SI{158}{rad/s} con
$V_2=\SI{6.9}{m/s}$ (dentro de C4) y trabado a \ang{-3} es el diseño base. La versión activa
con el servo permitiría la misma ley, pero exige un actuador por pala o un plato oscilante en
un cubo de \SI{44}{mm}, fuera del alcance del proyecto. La figura~\ref{fig:arranque-mec}(a)
muestra el mecanismo pasivo: el portapala gira sobre un buje en su eje longitudinal; una
masa de \SIrange{1}{2}{g} en un brazo excéntrico genera un momento de paso centrífugo
(efecto ``propeller moment'') que vence un resorte de torsión precargado.

\subsection{Acoplamiento \texorpdfstring{$\delta_3$}{delta-3}}\label{sec:arranque-d3}

El acoplamiento paso--batimiento es el método de mejor relación ganancia/costo. Con la bisagra
girada un ángulo $\delta_3$ en el plano del disco, el paso cambia con la conicidad,
\begin{equation}
  \theta=\theta_r-\tan\delta_3\,(\beta-\beta_r),
  \label{eq:arranque-d3}
\end{equation}
y como a bajas rpm la pala descansa en el tope de \ang{15} ($\beta-\beta_r\approx\ang{11}$), el
paso inicial resulta automáticamente \ang{-9} con $\delta_3=\ang{30}$ y \ang{-14} con
\ang{45}. A medida que la fuerza centrífuga baja la conicidad hacia \ang{4}, el paso vuelve
al de régimen: es un paso colectivo variable pasivo sin piezas añadidas. La simulación da
$t_{90}=\SI{0.76}{s}$ y $z_{90}=\SI{10.8}{m}$ con \ang{30} (\SI{-0.8}{s}, \SI{-10.9}{m}), y
\SI{0.71}{s} con \ang{45}. Las rpm finales cambian menos de \SI{1}{\percent} porque la
conicidad de régimen del modelo (\ang{4.5}) difiere algo de la de referencia; se corrige con
el paso del portapala.

Costos y riesgos: el acoplamiento también actúa en vuelo. Una ráfaga que suba la conicidad
\ang{2} baja el paso \ang{1.2} con \ang{30}, lo que sube las rpm y la fuerza centrífuga y
devuelve la pala: el efecto es estabilizante en batimiento pero modifica la frecuencia
$\nu_\beta$ y el amortiguamiento del capítulo~\ref{cap:estabilidad}, que habría que
recalcular. La geometría exige que la pala apoye en el tope sin torcer el pasador. El modelo
lineal~\eqref{eq:arranque-d3} se aplicó solo con la pala abierta ($\beta\ge0$).

\subsection{Flaps o pestañas de borde de fuga}

Una pestaña o flap de borde de fuga desviado cambia el ángulo de sustentación nula; se
modeló, en primer orden, como \ang{-4} de paso adicional en la mitad externa de la pala
(flap del \SI{20}{\percent} de cuerda desviado unos \ang{10}, con eficacia
$\approx0{,}4$, estimada). Fija, da $t_{90}=\SI{1.01}{s}$ pero lleva el régimen a
\SI{147}{rad/s} y $V_2=\SI{6.64}{m/s}$ (caso 9), igual que un paso de \ang{-6}: una pestaña
fija no es mejor que un paso más negativo. Solo tiene sentido si se retrae (por fuerza
centrífuga o presión dinámica), y entonces es un paso variable con un mecanismo más frágil
en la parte más solicitada de la pala. No se recomienda.

\subsection{Asimetría y arrancador tipo Savonius}

Un Savonius en el cubo ($\varnothing\SI{44}{mm}\times\SI{25}{mm}$, coeficiente de par
estático $C_{Q0}\approx0{,}3$, estimado) aporta
$Q=C_{Q0}\tfrac12\rho V^2 A\,r\approx\SI{0.8}{mN.m}$ a \SI{13.4}{m/s} y
\SI{0.2}{mN.m} a \SI{6.36}{m/s}: el \SI{3}{\percent} de $Q_0$, del orden del rozamiento. Sirve
como respaldo contra la trampa solo si el rozamiento es nominal y a alta $V$; además ocupa el
lugar de la copa del drogue. La asimetría entre palas (por ejemplo dos palas a \ang{-6} y dos
a \ang{0}) da el par promedio de su paso medio y añade desequilibrio aerodinámico y vibración
de 1P; no aporta nada que no dé el paso medio.

\subsection{Liberación a mayor velocidad}

El par crece con $V^2$, pero la inercia del rotor es la misma y la caída dura más: liberar a
\SI{16}{m/s} gana solo \SI{0.17}{s} y \SI{0.6}{m} respecto de \SI{13.4}{m/s}, y liberar a
\SI{6.36}{m/s} pierde \SI{0.56}{s} y \SI{1.7}{m} (casos 13 a 15). Con el drogue soltado antes
de abrir (v2), la velocidad de apertura sube por sí sola a \SIrange{14}{16}{m/s}. La altura casi no cambia
porque a mayor $V$ se cae más rápido durante el arranque. En la práctica la etapa 1 ya está
en su terminal (\SI{13.4}{m/s}) a la altura de liberación de C4, así que este método no tiene
recorrido.

\subsection{Reducir el rozamiento}

Con el paso de diseño el rozamiento cambia $t_{90}$ en \SI{\pm0.03}{s}: de 1,54 (aceite,
\SI{0.34}{mN.m}) a \SI{1.58}{s} (grasa fría, \SI{1.7}{mN.m}). Su valor es de margen: con
\SI{5}{mN.m} el paso crítico a \SI{6.36}{m/s} cae a \ang{-2.2} (tabla~\ref{tab:arranque-mapa}).
Por eso conviene rodamientos 688 con escudo metálico (ZZ, sin contacto) y no sellos de goma
(2RS), lavados y relubricados con aceite liviano o grasa de baja viscosidad, y un ensayo de
rozamiento a \SI{-10}{\celsius} si se esperan temperaturas bajas en altura.

%=======================================================================================
\section{Esquemas de los mecanismos}\label{sec:arranque-esquemas}

La figura~\ref{fig:arranque-mec} reúne los cuatro mecanismos con reparto de masas y
fuerzas. Las partes giratorias se dibujan en naranja, las fijas en azul y los elementos
elásticos o actuadores en rojo.

\begin{figure}[H]
\centering
\begin{subfigure}[b]{0.48\textwidth}\centering
\begin{tikzpicture}[scale=0.85]
  \fill[frot] (-0.5,-0.3) rectangle (3.6,0.3);
  \draw[rot] (-0.5,-0.3) rectangle (3.6,0.3);
  \node[font=\scriptsize] at (2.6,0) {portapala};
  \draw[eje] (-0.9,0) -- (4.1,0) node[right, font=\scriptsize] {eje de paso};
  \draw[rot, thick] (0.8,0.3) -- (0.8,1.3);
  \fill[cfreno] (0.8,1.45) circle (0.17);
  \node[font=\scriptsize, anchor=west] at (1.0,1.45) {contrapeso $m_c$};
  \draw[-{Stealth}, cfreno, thick] (0.8,1.65) -- (0.8,2.3) node[above, font=\scriptsize] {$m_c\Omega^2 r$};
  \draw[act, decorate, decoration={coil, aspect=0.6, segment length=4pt, amplitude=4pt}] (-0.4,0.3) -- (-0.4,1.4);
  \node[font=\scriptsize, anchor=east] at (-0.55,0.9) {resorte};
  \draw[-{Stealth}, crot, thick] (3.2,-0.6) arc (-60:60:0.7);
  \node[font=\scriptsize] at (3.2,-0.85) {$\theta$: $-8\to-3$};
\end{tikzpicture}
\caption{Paso variable centrífugo.}
\end{subfigure}\hfill
\begin{subfigure}[b]{0.48\textwidth}\centering
\begin{tikzpicture}[scale=0.85]
  \fill[ffijo] (-0.4,-1) rectangle (0.4,1);
  \draw[fijo] (-0.4,-1) rectangle (0.4,1);
  \node[font=\scriptsize, rotate=90] at (0,0) {cubo};
  \draw[rot, line width=2pt] (0.4,0) -- (4,0);
  \node[crot, font=\scriptsize, anchor=south] at (3,0.05) {pala};
  \draw[act, thick] (0.6,-0.9) -- (1.4,0.9);
  \node[cfreno, font=\scriptsize, anchor=west] at (1.4,0.9) {eje de bisagra};
  \draw[ref] (1,-1) -- (1,1);
  \draw[cgris] (1,0.6) arc (90:66:0.6);
  \node[font=\scriptsize] at (1.35,0.95) {};
  \node[font=\scriptsize] at (1.25,0.45) {$\delta_3$};
  \node[font=\scriptsize, align=center] at (2.6,-1.0) {$\beta\uparrow\Rightarrow\theta\downarrow$\\(vista en planta)};
\end{tikzpicture}
\caption{Bisagra con $\delta_3$.}
\end{subfigure}

\vspace{4mm}
\begin{subfigure}[b]{0.48\textwidth}\centering
\begin{tikzpicture}[scale=0.85]
  \fill[ffijo] (-0.2,-1.4) rectangle (0.2,1.4);
  \draw[fijo] (-0.2,-1.4) rectangle (0.2,1.4);
  \node[font=\scriptsize, anchor=east] at (-0.25,-1.2) {eje fijo};
  \draw[act, decorate, decoration={coil, aspect=0.5, segment length=3pt, amplitude=6pt}] (0,-0.6) -- (0,0.6);
  \fill[frot] (0.3,-0.9) rectangle (1.0,0.9);
  \draw[rot] (0.3,-0.9) rectangle (1.0,0.9);
  \fill[frot] (-1.0,-0.9) rectangle (-0.3,0.9);
  \draw[rot] (-1.0,-0.9) rectangle (-0.3,0.9);
  \node[font=\scriptsize, anchor=west] at (1.1,0.6) {cubo};
  \node[cfreno, font=\scriptsize, anchor=west] at (1.1,0) {resorte de torsión};
  \node[cfreno, font=\scriptsize, anchor=west] at (1.1,-0.4) {$d=\SI{1.1}{mm}$, 11 esp.};
  \draw[act] (1.0,-0.9) -- (1.4,-1.3) node[right, font=\scriptsize] {gatillo};
  \draw[-{Stealth}, cfreno, thick] (-0.9,-1.7) arc (200:340:1.0);
  \node[font=\scriptsize] at (0,-2.2) {reacción sobre el cuerpo};
\end{tikzpicture}
\caption{Resorte de pre-rotación en el cubo.}
\end{subfigure}\hfill
\begin{subfigure}[b]{0.48\textwidth}\centering
\begin{tikzpicture}[scale=0.85]
  \draw[rot, line width=1.5pt] (-1.8,0) .. controls (-0.6,0.25) and (0.6,0.25) .. (1.3,0.05);
  \draw[act, line width=1.5pt] (1.3,0.05) -- (1.95,-0.25);
  \draw[cgris, dashed] (1.3,0.05) -- (2.0,0.0);
  \node[font=\scriptsize] at (2.3,-0.05) {$\delta_f$};
  \node[cfreno, font=\scriptsize, anchor=north] at (1.6,-0.35) {flap/pestaña};
  \node[crot, font=\scriptsize] at (-0.4,0.55) {placa curvada};
  \draw[-{Stealth}, cfijo, thick] (-0.3,-1.3) -- (-0.3,-0.4);
  \node[cfijo, font=\scriptsize, anchor=west] at (-0.2,-1.0) {flujo de arranque};
\end{tikzpicture}
\caption{Pestaña de borde de fuga.}
\end{subfigure}
\caption{Esquemas (no a escala) de los mecanismos evaluados para inducir el arranque.}
\label{fig:arranque-mec}
\end{figure}

%=======================================================================================
\section{Matriz comparativa y recomendación}\label{sec:arranque-matriz}

La tabla~\ref{tab:arranque-matriz} compara los métodos respecto del caso base
($t_{90}=\SI{1.55}{s}$, $z_{90}=\SI{21.7}{m}$). La masa es una estimación; la complejidad y el
riesgo se califican de 1 (bajo) a 5 (alto). ``Compatible'' se refiere a los requisitos de
masa (\SI{500}{g}), envolvente (S2, S3), descenso (C4) y video estable.

\begin{table}[H]
\centering
\small
\caption{Matriz comparativa de los métodos de arranque.}
\label{tab:arranque-matriz}
\begin{tabularx}{\textwidth}{lrrrccY}
\toprule
Método & $\Delta t_{90}$ [s] & $\Delta z_{90}$ [m] & Masa [g] & Compl. & Riesgo & Compatibilidad / comentario \\
\midrule
Paso \ang{-4} (fijo) & $-0{,}26$ & $-3{,}6$ & 0 & 1 & 1 & Sí. $V_2$ $+\SI{0.07}{m/s}$, +\SI{5}{\percent} rpm \\
Paso \ang{-6} (fijo) & $-0{,}56$ & $-7{,}6$ & 0 & 1 & 2 & Sí. $V_2=6{,}60$; más carga y flameo \\
Torsión $+\ang{10}$ & $-0{,}15$ & $-2{,}1$ & 0 & 3 & 2 & Sí. +\ang{0.6} de margen; molde \\
Resorte $\Omega_0=40$ & $-0{,}94$ & $-13{,}8$ & 8--10 & 4 & 4 & Gira el cuerpo $\approx\SI{76}{rad/s}$ \\
Servo (trinquete) & $\approx-0{,}3$ & $\approx-5$ & 2 & 3 & 3 & Servo ocupado en la traba; reacción \\
Motor \SI{3}{mN.m}, \SI{1}{s} & $-0{,}07$ & $-1{,}3$ & 4--6 & 4 & 3 & Consumo, cableado al cubo, reacción \\
Apertura con bisagra girada & $\approx-0{,}2$ & $\approx-3$ & 0 & 2 & 2 & Es un $\delta_3$ (fila siguiente) \\
Paso variable centrífugo & $-0{,}60$ & $-8{,}3$ & 8--12 & 5 & 3 & Falla benigna; cubo \SI{44}{mm} saturado \\
$\delta_3=\ang{30}$ & $-0{,}79$ & $-10{,}9$ & 0 & 2 & 3 & Recalcular $\nu_\beta$ y tope \\
Pestaña fija (\ang{-4} eq.) & $-0{,}54$ & $-7{,}4$ & 1 & 2 & 2 & Igual a paso \ang{-6}; retráctil: no \\
Savonius en el cubo & $\approx-0{,}02$ & $\approx0$ & 3 & 2 & 2 & Choca con la copa del drogue \\
Liberación a \SI{16}{m/s} & $-0{,}17$ & $-0{,}6$ & 0 & --- & --- & Se da sola al soltar el drogue \\
Rozamiento $0{,}73\to0{,}34$ & $-0{,}01$ & $-0{,}1$ & 0 & 1 & 1 & Sí. Gana margen a baja $V$ \\
\bottomrule
\end{tabularx}
\end{table}

\begin{resultado}[Recomendación]
\begin{enumerate}
  \item \textbf{Paso fijo entre \ang{-3} y \ang{-4}, sin mecanismo de arranque.} El arranque es
        seguro (margen de \ang{2.7} a \ang{3.9} sobre el paso crítico a 6,36--\SI{13.4}{m/s},
        más que la tolerancia de $\pm\ang{1}$) y la caída de arranque, 18--\SI{22}{m}, cabe en
        la altura de la etapa 2. Pasar a \ang{-4} gana \SI{3.6}{m} con
        $\Delta V_2=+\SI{0.07}{m/s}$: es la mejora gratis y se recomienda si el portapala
        se reimprime.
  \item \textbf{Rodamientos ZZ limpios y lubricados con aceite liviano}; prohibir sellos de
        contacto. Cuesta nada y protege el margen a baja $V$, donde
        $Q_f=\SI{5}{mN.m}$ lo agotaría.
  \item \textbf{Si se necesita altura}, el siguiente paso es la bisagra con
        $\delta_3\approx\ang{30}$: \SI{-11}{m} sin masa añadida, mejor que el resorte
        (\SI{-14}{m} pero con \SI{76}{rad/s} de giro del cuerpo y \SI{10}{g}) y que el paso
        variable centrífugo (\SI{-8}{m}, mecanismo complejo). Requiere recalcular el
        batimiento (cap.~\ref{cap:estabilidad}) y validarlo en banco.
  \item \textbf{Descartados}: motor, servo, Savonius, pestaña retráctil y liberación más
        temprana, por ganancia menor que la incertidumbre del modelo
        (\SI{0.18}{s}) o por la reacción sobre el cuerpo.
\end{enumerate}
\end{resultado}

%=======================================================================================
\section{Plan de ensayo del arranque en banco}\label{sec:arranque-ensayo}

El ensayo debe medir las dos magnitudes de las que depende el arranque, $Q_f$ y $Q_0$, y
luego verificar el transitorio. Ninguno se hizo todavía.

\begin{table}[H]
\centering
\small
\caption{Ensayos de arranque y criterios de aceptación.}
\label{tab:arranque-ensayo}
\begin{tabularx}{\textwidth}{lYY}
\toprule
Ensayo & Método & Criterio \\
\midrule
A1. Rozamiento & Cubo con palas montado en el eje; impulso a mano a $\approx\SI{100}{rad/s}$; desaceleración por sensor hall; $Q_f=I_z|\dot\Omega|$ a baja $\Omega$ (sin palas para separar el arrastre). Ambiente y \SI{-10}{\celsius}. & $Q_f\le\SI{1.7}{mN.m}$ a \SI{20}{\celsius}; $\le\SI{3}{mN.m}$ en frío. \\
A2. Par estático & Rotor trabado con brazo y balanza (o celda de carga de \SI{100}{g}) frente a un ventilador o túnel con flujo vertical de 5, 7, 10 y \SI{13}{m/s} (anemómetro en el plano del disco). & $Q_0\ge\SI{15}{mN.m}$ a \SI{13}{m/s} con \ang{-3} (la mitad del modelo); signo positivo a todo $V$. \\
A3. Barrido de paso & Repetir A2 con portapalas de \ang{-1}, \ang{0} y \ang{+1}. & Ubicar el paso crítico; aceptar si $\theta_c\ge\ang{-1.5}$ (margen $\ge\ang{1.5}$). \\
A4. Arranque libre & Rotor libre en el flujo, palas colgando y trabadas; liberación por el servo real; registro de rpm a \SI{1}{kHz} y video a 240 fps. 20 repeticiones. & 20/20 arranques; $t_{90}$ en banco dentro de $\pm\SI{40}{\percent}$ del modelo corrido a la $V$ del banco; apertura $<\SI{0.3}{s}$. \\
A5. Trampa & Con \ang{+1}: confirmar que no arranca (la caja cae a $\approx\SI{14}{m/s}$); con impulso a mano medir el $\Omega_0$ umbral. & Umbral dentro de $\pm\SI{50}{\percent}$ de la tabla~\ref{tab:arranque-om0}; valida el modelo de la rama lenta. \\
A6. Caída real & Lanzamiento desde dron a $\ge\SI{60}{m}$ con la etapa 1 completa. & $z_{90}\le\SI{35}{m}$; $V$ final dentro de C4. \\
\bottomrule
\end{tabularx}
\end{table}

\begin{nota}[Límites del banco]
Un ventilador da un chorro de \SI{30}{cm} a \SI{1}{m} con perfil no uniforme y turbulencia
alta, y el rotor de \SI{40}{cm} lo bloquea: las $V$ medidas son de referencia, no
absolutas. Por eso los criterios de A2 y A4 son relativos al modelo y amplios, y la
validación final es el ensayo A6. Si A3 da un paso crítico por debajo de \ang{-2}, el diseño
pasa a \ang{-4} sin más análisis.
\end{nota}

\paragraph{Límites del modelo.} Las predicciones de este capítulo heredan las del BEMT en el
rango $a>1$ (rotor casi parado), la polar en pérdida profunda con su incertidumbre de
$C_{D,\max}$ entre 1,17 y 1,98, un batimiento cuasiestático después de la apertura y un
rozamiento constante. No incluyen la inercia del flujo inducido, que retrasa la respuesta
del par en unas pocas revoluciones (del orden de $R/v_i\approx\SI{0.05}{s}$) y es menor que
las demás incertidumbres. Los mecanismos ($\delta_3$, paso variable, resorte) se modelaron con
leyes ideales; sus fricciones y holguras propias no se consideraron.
